\documentclass[UTF8,a4paper,12pt]{ctexart}
\usepackage{ctex}
\usepackage{amsmath}
\numberwithin{equation}{section}
\allowdisplaybreaks[4]       %多行公式中换页
\usepackage{array}
\usepackage[font=small,font=bf,labelsep=none]{caption}
\usepackage{amssymb}
\usepackage{tikz}
\usepackage{amsthm}
\usepackage{mathrsfs}
\usepackage{dutchcal}
\usepackage{color}
\usepackage{graphicx}    %插入图片
\usepackage{times}
\usepackage{mathptmx}
\usepackage{fancyhdr} %页眉页脚
\usepackage[backend=biber, style=gb7714-2015]{biblatex}
\usepackage{bicaption}
\usepackage{float}  % 导言区引入
\usepackage[colorlinks=true, linkcolor=blue, urlcolor=blue]{hyperref}
% 文献表条目间的间距
\setlength{\bibitemsep}{0pt}
% 导入参考文献数据库
\addbibresource{main.bib}




\pagestyle{fancy}
\fancyhf{}
\fancyfoot[C]{\thepage}
\usepackage{setspace}
\setlength{\baselineskip}{20pt}
\newcommand*{\circled}[1]{\lower.7ex\hbox{\tikz\draw (0pt, 0pt)%
    circle (.5em) node {\makebox[1em][c]{\small #1}};}}
\usepackage{hyperref}  %目录
\hypersetup{colorlinks=true,linkcolor=black}
\captionsetup[figure][bi-second]{name=Figure} %设置图的英文编号前缀
\captionsetup[table][bi-second]{name=Table} %设置表的英文编号前缀
\numberwithin{equation}{section}%公式按章节编号
\numberwithin{figure}{section}%图表按章节编号
\numberwithin{table}{section}
\renewcommand {\thefigure} {\thesection{}-\arabic{figure}}%设定图片的编号。这样设置的实现效果为图1-1
\renewcommand {\thetable} {\thesection{}-\arabic{figure}}
\usepackage[justification=centering]{caption} %保持居中，防止标题过长导致左对齐
\usepackage{caption}
\captionsetup{font={small},labelsep=quad}%文字5号，之间空一个汉字符位。
\captionsetup[table]{font={bf}} %表格表号与表题加粗
\usepackage{appendix}
\usepackage{tocloft} 
\renewcommand{\cftsecleader}{\cftdotfill{\cftdotsep}} %为目录中section补上引导点
\usepackage{titletoc}
\titlecontents{section}[0pt]{\addvspace{6pt}\filright\bf}%
               {\contentspush{\thecontentslabel \quad }}%
               {}{\titlerule*[8pt]{.}\contentspage}
\makeatletter %双线页眉
\def\headrule{{\if@fancyplain\let\headrulewidth\plainheadrulewidth\fi%
\hrule\@height 1.5pt \@width\headwidth\vskip1.5pt%上面线为1pt粗
\hrule\@height 0.5pt\@width\headwidth  %下面0.5pt粗
\vskip-2\headrulewidth\vskip-1pt}      %两条线的距离1pt
  \vspace{6mm}}     %双线与下面正文之间的垂直间距
\makeatother
\CTEXsetup[format={\heiti \zihao{3} \bfseries \center}]{section}
\CTEXsetup[number={实验\chinese{section}}]{section} 

% 自定义章节目录格式
% \ctexset{section={
%   format={\heiti \zihao{3} \bfseries \center},
%   number={第\chinese{section}章}
% }}
\usepackage[explicit]{titlesec}
\titlespacing*{\section}{0pt}{24pt plus .24pt minus .24pt}{18pt plus .0ex}




\begin{document}

\zihao{-4}\songti

\thispagestyle{empty}

\renewcommand{\headrulewidth}{0pt}
\begin{figure}[htb]
\centering % 确保图片和文字居中
\includegraphics[width=1\textwidth]{fig/北京理工大学校徽.png} % 设置图片宽度为页面宽度的80%
\end{figure}

~\\
%该页为中文扉页。无需页眉页脚，纸质论文应装订在右侧
~\\
\begin{center}
\songti \zihao{2} \textbf{XXX实验课程报告}
\end{center}
%中文论文标题，1行或2行，宋体，加粗，二号，居中。论文题目不得超过36个汉字
~\\
~\\
~\\
\begin{center}
\heiti \zihao{3}
\begin{tabular}{l}



\textbf{姓\qquad 名：}XXX\\
\textbf{学\qquad 号：}112022XXXX\\
\textbf{班\qquad 级：}0611XXXX\\ 
\textbf{学\qquad 院：}自动化学院\\ 
\textbf{专\qquad 业：}自动化\\ 
\textbf{指导老师：}XX\\

\end{tabular}
\end{center}
~\\
\begin{center}
\songti \zihao{4} \textbf{\today}
\end{center}

\newpage

\pagenumbering{Roman}

\pagenumbering{arabic}

\fancyhead[LH]{传感器综合实验课程报告}
\fancyhead[RH]{实验一\quad 电阻应变片式传感器实验}
\section{电阻应变片式传感器实验}

\begin{center}
    {\zihao{-4}
        小组成员：\underline{\quad 闵鹏远、王安彤 \quad} \quad
        实验时间：\underline{\quad 2025年3月28日 \quad}
    }
\end{center}

\subsection{应变片双臂特性试验}

\subsubsection{实验目的}

\begin{enumerate}
    \item 了解金属箔式应变片的工作原理与应变效应
    \item 掌握应变片式电阻传感器的
    \item 掌握机械回程差的消除方法。
\end{enumerate}

\subsubsection{实验内容}

针对金属箔式应变片，构建半桥电桥电路，如下图，将电阻应变式传感器的电阻变化转换成电压或电流信号。通过调节测微头微分筒改变梁的受力程度，测试电阻应变片的特性。

\subsubsection{实验步骤}

\subsubsection{实验结果}

1. 原始数据

2. 曲线图

3. 灵敏度

4. 非线性误差

\subsection{应变片温度特性试验及温度补偿实验}

\subsubsection{实验目的}

\subsubsection{实验内容}

\subsubsection{实验步骤}

\subsubsection{实验结果}

1. 原始数据

2. 计算

3. 分析比较

\subsection{应变片直流全桥的应用——电子秤实验}

\subsubsection{实验记录}

\subsubsection{电压与重量的关系公式}

\subsubsection{实验结果分析}

\subsubsection{称重两个物品重量的计算}

\subsection{测试样例}

这里用于测试内容，请后续直接删除即可 

图片引用测试\ref{fig:myfig2}：

\begin{figure}[H] % 这里 H 表示 强制放置在当前位置，更改为htbp表示根据需要放置
    \centering
    \includegraphics[width=0.8\textwidth]{fig/北京理工大学校徽.png} % 可调整宽度
    \caption{测试图片}
    \label{fig:myfig2}
\end{figure}

公式示例\ref{eq:pid-incremental}：
\begin{equation}
    \left\{
    \begin{aligned}
        u(k) &= q_0 e(k) + B(k-1) \\
        B(k) &= u(k) + q_1 e(k-1) + q_2 e(k-2) \\
        q_0 &= K_p \left(1 + \frac{T_s}{T_i} + \frac{T_d}{T_s} \right) \\
        q_1 &= -K_p \left(1 + 2 \cdot \frac{T_d}{T_s} \right) \\
        q_2 &= K_p \cdot \frac{T_d}{T_s}
    \end{aligned}
    \right.
    \label{eq:pid-incremental}
    \end{equation} 

表格示例：
\begin{table}[H]
    \centering
    \caption{高频感应加热的基本参数}
    \begin{tabular}{|c| c|c|c|}
    \hline
    感应频率 &感应发生器功率 & 工件移动速度  &感应圈与零件间隙\\
    （KHz）&($\% \times$80Kw) &(mm/min)  &(mm)\\
    \hline
    250 &88 &5900 &1.65\\
    \hline
    250 &88 &5900 &1.65\\
    \hline
    250 &88 &5900 &1.65\\
    \hline
    250 &88 &5900 &1.65\\
    \hline
    \end{tabular}
\end{table}


\begin{table}[H]
    \centering
    \captionsetup{singlelinecheck=off}
    \caption*{续表} %取消编号
    \begin{tabular}{|c| c|c|c|}
    \hline
    感应频率 &感应发生器功率 & 工件移动速度  &感应圈与零件间隙\\
    （KHz）&($\% \times$80Kw) &(mm/min)  &(mm)\\
    \hline
    250 &88 &5900 &1.65\\
    \hline
    250 &88 &5900 &1.65\\
    \hline
    \end{tabular}
\end{table}
%表格太大需要转页时，需要在续表上方注明“续表”，表头也应重复排出。 % 这里的 tab/tab1.tex 是表格文件的路径

\newpage

\fancyhead[LH]{传感器综合实验课程报告}
\fancyhead[RH]{实验二\quad 其他典型传感器实验}
\section{其他典型传感器实验}

\begin{center}
    {\zihao{-4}
        小组成员：\underline{\quad 闵鹏远、王安彤 \quad} \quad
        实验时间：\underline{\quad 2025年4月6日 \quad}
    }
\end{center}

\subsection{压阻式压力传感器实验}

\subsubsection{实验目的}

\subsubsection{实验内容}

\subsubsection{实验步骤}

\subsubsection{实验结果}

\subsection{电容式传感器静态特性试验}

\subsubsection{实验目的}

\subsubsection{实验内容}

\subsubsection{实验步骤}

\subsubsection{实验结果}

% 气敏传感器和湿敏传感器实验报告没有强制要求撰写，视自身情况而定

% \subsection{气敏传感器实验}

% \subsubsection{实验目的}

% \subsubsection{实验内容}

% \subsubsection{实验步骤}

% \subsubsection{实验结果}

% \subsection{湿敏传感器}

% \subsubsection{实验目的}

% \subsubsection{实验内容}

% \subsubsection{实验步骤}

% \subsubsection{实验结果}

\subsection{光纤位移传感器测位移特性试验}

\subsubsection{实验目的}

\subsubsection{实验内容}

\subsubsection{实验步骤}

\subsubsection{实验结果}

\newpage

\fancyhead[LH]{传感器综合实验课程报告}
\fancyhead[RH]{实验三\quad 霍尔传感器实验}
\section{霍尔传感器实验}

\begin{center}
    {\zihao{-4}
        小组成员：\underline{\quad 闵鹏远、王安彤 \quad} \quad
        实验时间：\underline{\quad 2025年4月9日 \quad}
    }
\end{center}

\subsection{线性霍尔式传感器位移特性实验}

\subsubsection{实验目的}

\subsubsection{实验内容}

\subsubsection{实验步骤}

\subsubsection{实验结果}

\subsection{基于霍尔传感器的电子秤设计}

\subsubsection{设计目的}

\subsubsection{设计任务}

\subsubsection{设计要求}

\subsubsection{基于霍尔式传感器的电子秤原理及原理图}

\subsubsection{选用实验部件及器件}

\subsubsection{硬件电路设计步骤及操作步骤}

\subsubsection{实验结果}

1. 电压与重量的关系记录

2. 电压-重量曲线

3. 计算线性范围及灵敏度

4. 计算称重物品的重量

\subsubsection{总结分析}

\newpage

\fancyhead[LH]{传感器综合实验课程报告}
\fancyhead[RH]{实验四\quad 基于光敏电阻的设计实验}
\section{基于光敏电阻的设计实验}

\begin{center}
    {\zihao{-4}
        小组成员：\underline{\quad XXX、XXX \quad} \quad
        实验时间：\underline{\quad 20XX年X月X日 \quad}
    }
\end{center}
\subsection{光敏电阻的暗光亮灯实验}

\subsubsection{实验目的}

\subsubsection{实验元件}

\subsubsection{实验原理}

\subsubsection{实验步骤}

\subsubsection{实验结果}

\subsection{光敏电阻闪灯电路实验设计}

\subsubsection{设计目的}

\subsubsection{设计要求}

\subsubsection{设计原理图}

\subsubsection{选用实验部件及器件}

\subsubsection{实验结果}

\newpage

\fancyhead[LH]{传感器综合实验课程报告}
\fancyhead[RH]{实验五\quad 光电传感器实验}
\section{光电传感器实验}

\begin{center}
    {\zihao{-4}
        小组成员：\underline{\quad XXX、XXX \quad} \quad
        实验时间：\underline{\quad 20XX年X月X日 \quad}
    }
\end{center}

\subsection{XXX}

\newpage

\fancyhead[LH]{传感器综合实验课程报告}
\fancyhead[RH]{实验六\quad 温度变送器虚拟仿真实验}
\section{温度变送器虚拟仿真实验}

\begin{center}
    {\zihao{-4}
        小组成员：\underline{\quad 闵鹏远、王安彤 \quad} \quad
        实验时间：\underline{\quad 2025年4月XX日 \quad}
    }
\end{center}

\subsection{实验目的}

\subsection{实验原理}

\subsection{实验内容及步骤}

\subsection{实验记录表格}
仿真实验系统中的截图，包括实验数据表格及成绩页面

\subsection{结果分析}

\newpage

\end{document} 